\section{Backbone Architecture}

The model employs a pre-trained ResNet-34 as its encoder backbone,
leveraging its deep residual learning framework for efficient feature
extraction. The choice of ResNet-34 provides several key advantages:
efficient feature extraction through residual connections,
pre-trained weights that capture rich visual representations, and
stable gradient flow during training \cite{HeEtAl2015}. The
ResNet-34 backbone is modified to serve as the encoder in our UNET
architecture by removing the final fully connected layer and
utilizing the feature maps from different stages of the network for
skip connections.

\subsection{U-Net Architecture}

The U-Net architecture follows a traditional encoder-decoder structure
with skip connections, where the encoder path implements the
ResNet-34 structure. The decoder path employs transposed convolutions
for upsampling, creating a symmetrical architecture that effectively
captures both high-level and low-level features. The architecture
incorporates four downsampling stages in the encoder, corresponding
to the ResNet-34 blocks, and four upsampling stages in the decoder.
These stages are connected through skip connections that bridge
corresponding encoder and decoder stages, allowing the network to
preserve fine-grained details. Each convolution operation is followed
by batch normalization and ReLU activation to ensure stable training
and effective feature learning.

This U-Net architecture is expected to be coupled with a reliability
branch which leads to a multi-output network. Three different approaches
are considered for this branch: \textbf{scalar}, \textbf{features},
and \textbf{pixel}.
